\chapter[Band 13: Schranken, Infima und Suprema]{Band 13 auf einen Blick}
\noindent\textbf{Schranken, Infima und Suprema}\par\medskip
\noindent\emph{Lesart.} Es gilt \(\PartOrd{A,\leq}\),
\(T\subseteq A\). In dieser Übersicht kürzen wir die festen
Trägerindizes ab: \(L(T)=\LB_{A,\leq}(T)\),
\(U(T)=\UB_{A,\leq}(T)\), \(\inf T=\inf_{A,\leq}(T)\),
\(\sup T=\sup_{A,\leq}(T)\). Infimum und Supremum werden nur
unter ihrer jeweiligen Existenzvoraussetzung verwendet.
\begin{BandUebersicht}
\textbf{Untere und obere Schranken}
\UebFormel{\begin{aligned}
L(T)&\coloneqq\{a\in A\mid\forall x\in T\,(a\leq x)\},\\
U(T)&\coloneqq\{a\in A\mid\forall x\in T\,(x\leq a)\}.
\end{aligned}}
\UebQuellen{\FormulaRefAuto{LowerBoundSet}[def];
\FormulaRefAuto{UpperBoundSet}[def]}
&
\textit{Elementkriterien und Typisierung}
\UebFormel{\begin{aligned}
&u\in L(T)\leftrightarrow u\in A\land\forall x\in T\,(u\leq x),\\
&u\in U(T)\leftrightarrow u\in A\land\forall x\in T\,(x\leq u),\\
&L(T)\subseteq A,\qquad U(T)\subseteq A.
\end{aligned}}
\UebQuellen{\FormulaRefAuto{LowerBoundSetCharacterization}[thm];
\FormulaRefAuto{UpperBoundSetCharacterization}[thm];
\FormulaRefAuto{LowerBoundSetSubsetCarrier}[thm];
\FormulaRefAuto{UpperBoundSetSubsetCarrier}[thm]}\\
\addlinespace[.8ex]
\textbf{Schranken eines Elementpaares}
Die allgemeinen Schrankenmengen werden auf \(T=\{x,y\}\)
mit \(x,y\in A\) angewandt.
&
\textit{Zwei Ungleichungen genügen}
\UebFormel{\begin{aligned}
&u\in U(\{x,y\})\\
&\quad\leftrightarrow u\in A\land x\leq u\land y\leq u,\\
&d\in L(\{x,y\})\\
&\quad\leftrightarrow d\in A\land d\leq x\land d\leq y.
\end{aligned}}
\UebQuellen{\FormulaRefAuto{PairUpperBoundSetCharacterization}[thm];
\FormulaRefAuto{PairLowerBoundSetCharacterization}[thm]}\\
\addlinespace[.8ex]
\textbf{Supremum}
Existiert eine kleinste obere Schranke, setzt man
\UebFormel{\sup T\coloneqq\min_{A,\leq}(U(T)).}
Die ausdrückliche Voraussetzung ist
\UebFormel{\exists s\in U(T)\,\forall u\in U(T)\,(s\leq u).}
\UebQuellen{\FormulaRefAuto{Supremum}[def]}
&
\textit{Kleinste obere Schranke}
Unter der linken Existenzvoraussetzung:
\UebFormel{\begin{aligned}
&\sup T\in A,\qquad\sup T\in U(T),\\
&x\in T\vdash x\leq\sup T,\\
&u\in U(T)\vdash\sup T\leq u.
\end{aligned}}
\UebQuellen{\FormulaRefAuto{SupremumInCarrier}[thm];
\FormulaRefAuto{SupremumBelongsUpperBounds}[thm];
\FormulaRefAuto{SupremumDominatesMember}[thm];
\FormulaRefAuto{SupremumLeastUpperBound}[thm]}\\
\addlinespace[.8ex]
\textbf{Infimum}
Existiert eine größte untere Schranke, setzt man
\UebFormel{\inf T\coloneqq\max_{A,\leq}(L(T)).}
Die ausdrückliche Voraussetzung ist
\UebFormel{\exists i\in L(T)\,\forall d\in L(T)\,(d\leq i).}
\UebQuellen{\FormulaRefAuto{Infimum}[def]}
&
\textit{Größte untere Schranke}
Unter der linken Existenzvoraussetzung:
\UebFormel{\begin{aligned}
&\inf T\in A,\qquad\inf T\in L(T),\\
&x\in T\vdash\inf T\leq x,\\
&d\in L(T)\vdash d\leq\inf T.
\end{aligned}}
\UebQuellen{\FormulaRefAuto{InfimumInCarrier}[thm];
\FormulaRefAuto{InfimumBelongsLowerBounds}[thm];
\FormulaRefAuto{InfimumBelowMember}[thm];
\FormulaRefAuto{InfimumGreatestLowerBound}[thm]}\\
\addlinespace[.8ex]
\textbf{Nachweis eines Kandidaten}
Um einen vorgegebenen Wert als Supremum oder Infimum zu
identifizieren, prüft man Schranken- und Extremaleigenschaft.
&
\textit{Eindeutige Identifikation}
\UebFormel{\begin{aligned}
&u\in U(T)\dsep\forall a\in U(T)\,(u\leq a)\\
&\hspace{12mm}\vdash u=\sup T,\\
&u\in L(T)\dsep\forall a\in L(T)\,(a\leq u)\\
&\hspace{12mm}\vdash u=\inf T.
\end{aligned}}
\UebQuellen{\FormulaRefAuto{SupremumCandidateEqualsSupremum}[thm];
\FormulaRefAuto{InfimumCandidateEqualsInfimum}[thm]}\\
\addlinespace[.8ex]
\textbf{Duale Ordnung}
\UebFormel{x\geq y\leftrightarrow y\leq x\quad(x,y\in A).}
Beim Umkehren der Ordnung wechseln untere und obere Schranken ihre Rolle.
\UebQuellen{\FormulaRefAuto{OrderDualRelation}[def]}
&
\textit{Dualität der Extremalbegriffe}
\UebFormel{\begin{aligned}
&\LB_{A,\geq}(T)=\UB_{A,\leq}(T),\\
&\UB_{A,\geq}(T)=\LB_{A,\leq}(T),\\
&\sup_{A,\leq}(T)=\inf_{A,\geq}(T),\\
&\inf_{A,\leq}(T)=\sup_{A,\geq}(T).
\end{aligned}}
Die letzten beiden Gleichungen setzen die Existenz des links
stehenden Extremalwertes voraus.
\UebQuellen{\FormulaRefAuto{DualLowerBoundsAreUpperBounds}[thm];
\FormulaRefAuto{DualUpperBoundsAreLowerBounds}[thm];
\FormulaRefAuto{DualSupremumEqualsInfimum}[thm];
\FormulaRefAuto{DualInfimumEqualsSupremum}[thm]}\\
\end{BandUebersicht}
